\chapter{اللقاء السابع والثلاثون: تفضيل الرسل وآية الكرسي وبداية (لا إكراه في الدين)}

\section{مقدمة}
السلام عليكم ورحمة الله وبركاته. الحمد لله رب العالمين، وأشهد أن لا إله إلا الله وحده لا شريك له، وأشهد أن محمداً عبده ورسوله، صلى الله عليه وعلى آله وصحبه وسلم.

استأنف الشيخ هذا اللقاء من قوله تعالى: \begin{ayah}تِلْكَ الرُّسُلُ فَضَّلْنَا بَعْضَهُمْ عَلَى بَعْضٍ\end{ayah}، ثم انتقل إلى آية الإنفاق قبل مجيء يوم القيامة، ووقف طويلاً مع آية الكرسي، ثم افتتح باب قوله تعالى: \begin{ayah}لَا إِكْرَاهَ فِي الدِّينِ\end{ayah} وما بعده من قوله: \begin{ayah}اللَّهُ وَلِيُّ الَّذِينَ آمَنُوا\end{ayah}، وختم ببداية قصة محاجة إبراهيم عليه السلام.

وهذا اللقاء من اللقاءات الجامعة جداً؛ لأنه يبين أصولاً في تفضيل الأنبياء، وسنن الاختلاف بعد مجيء البينات، وفقه الإنفاق قبل فوات زمن العمل، ويبرز منهج \rawi{الطبري} في أبواب الصفات، والعموم والخصوص، والنسخ، والاختلاف الذي يرد إلى المعنى الجامع لا إلى التضاد.

\separator

\section{تفضيل الرسل وبعض وجوه الاصطفاء}
\isnad{قال أبو جعفر:} المراد بقوله تعالى: \begin{ayah}تِلْكَ الرُّسُلُ\end{ayah} الرسل الذين قص الله أخبارهم في هذه السورة، كإبراهيم وموسى وعيسى وداود وغيرهم.

وبيّن \rawi{الطبري} أن من وجوه هذا التفضيل:
\begin{itemize}
    \item أن الله كلم موسى عليه السلام.
    \item وأنه رفع بعضهم درجات.
    \item وأن من أعظم وجوه ذلك إرسال محمد صلى الله عليه وسلم إلى الناس كافة.
\end{itemize}

ونبّه الشيخ إلى معنى عظيم هنا: أن الاصطفاء والاجتباء والاختصاص من الله منّة تستوجب مزيد شكر وقيام بحق الله، لا مجرد شعور بالفضل. فكلما خص الله عبداً بنعمة أو علم أو فتح، كان ذلك أدعى إلى الخوف من التقصير، لا إلى الركون إلى المنزلة.

\begin{fawaid}[كل اصطفاء من الله يستدعي شكراً يوازيه]
المنح الإلهية ليست مجرد تشريف، بل هي كذلك تكليف وسؤال؛ فمن اصطفاه الله بعلم أو دعوة أو عبادة أو منزلة، وجب عليه أن يكون أشد الناس شكراً وقياماً بحق تلك النعمة.
\end{fawaid}

\section{عيسى عليه السلام والبينات وروح القدس}
لما قال الله تعالى: \begin{ayah}وَآتَيْنَا عِيسَى ابْنَ مَرْيَمَ الْبَيِّنَاتِ وَأَيَّدْنَاهُ بِرُوحِ الْقُدُسِ\end{ayah}، أجمل \rawi{الطبري} هنا ما سبق بسطه في مواضع قبلها؛ لأن من طريقته أنه إذا توسع في أصل من الأصول أو لفظ من الألفاظ، لم يعد تكراره كل مرة، بل يحيل على ما سبق.

ومن الفوائد المنهجية التي أبرزها الشيخ هنا أن سورة البقرة مثال عظيم على منهج \rawi{الطبري} كله تقريباً، ولذلك كانت أوسع مواضع الكتاب في تقرير أصوله في اللغة، والآثار، والفقه، وأبواب العقيدة، والنسخ، والعموم والخصوص.

\begin{fawaid}[سورة البقرة أنموذج مركز لمنهج الطبري]
من أراد أن يفهم طريقة الطبري في التفسير فليعتن بسورة البقرة عناية خاصة؛ لأنها جمعت أكثر الأبواب التي يبني عليها تفسيره من جهة اللغة، والآثار، والفقه، والاعتقاد، ومناهج الترجيح.
\end{fawaid}

\section{وجود البينات لا يمنع وقوع الاختلاف}
في قوله تعالى: \begin{ayah}وَلَوْ شَاءَ اللَّهُ مَا اقْتَتَلَ الَّذِينَ مِنْ بَعْدِهِمْ مِنْ بَعْدِ مَا جَاءَتْهُمُ الْبَيِّنَاتُ\end{ayah}، قرر \rawi{الطبري} أن اختلاف الناس واقتتالهم بعد مجيء الرسل والبينات إنما وقع لما شاء الله كونه وقدره، لا لأن الحجة لم تقم عليهم.

ووقف الشيخ هنا مع تربية مهمة للدعاة وطلبة العلم: أن وجود البيان والدليل لا يستلزم أن ينقاد كل أحد، ولا أن ينتهي الخلاف في الناس. فربما قامت البينة أوضح ما تكون، ثم لا يزيد ذلك بعض القلوب إلا عناداً. ولذلك كان من تمام فقه الداعية أن يعلم ما كلف به وما لم يكلف به:
\begin{itemize}
    \item فكلف بالبلاغ والبيان.
    \item ولم يكلف بإدخال الإيمان في القلوب.
\end{itemize}

\begin{fawaid}[على الداعية أن يفرّق بين إقامة الحجة وبين حصول الاستجابة]
قيام البينة واجب على أهل العلم والدعوة، أما حصول الإيمان في القلوب فبيد الله. ومن لم يفرق بين الأمرين أتعب نفسه بما لم يكلفه الله به.
\end{fawaid}

\section{الإرادة الكونية والإرادة الشرعية}
استخرج الشيخ من هذا الموضع فرقاً نافعاً بين نوعي الإرادة في القرآن:
\begin{itemize}
    \item إرادة شرعية: وهي ما يحبه الله ويرضاه ويأمر به، وقد يقع وقد لا يقع.
    \item وإرادة كونية قدرية: وهي ما شاء الله وقوعه في خلقه، ولا بد أن يقع.
\end{itemize}

ففي مثل قوله تعالى: \begin{ayah}وَلَوْ شَاءَ اللَّهُ مَا اقْتَتَلُوا\end{ayah}، المراد مشيئته الكونية القدرية. أما في نحو: \begin{ayah}وَاللَّهُ يُرِيدُ أَنْ يَتُوبَ عَلَيْكُمْ\end{ayah}، فالمراد الإرادة الشرعية.

\begin{fawaid}[التمييز بين الإرادتين يرفع إشكالات كثيرة في التفسير]
إذا لم يميز القارئ بين الإرادة الشرعية التي تتعلق بالمحبة والأمر، والإرادة الكونية التي تتعلق بالمشيئة والوقوع، اختلطت عليه آيات كثيرة في القرآن.
\end{fawaid}

\separator

\section{الإنفاق قبل مجيء يوم لا بيع فيه ولا خلة ولا شفاعة}
لما جاء قوله تعالى: \begin{ayah}يَا أَيُّهَا الَّذِينَ آمَنُوا أَنْفِقُوا مِمَّا رَزَقْنَاكُمْ\end{ayah}، بيّن \rawi{الطبري} أن أول ما يدخل في هذه النفقة أداء الحقوق الواجبة: كالزكاة، والنفقات الواجبة، وسائر الحقوق المتعلقة بالمال، لا أن يحمل الخطاب ابتداءً على التطوع مع تضييع الفرائض.

ثم شرح معنى قوله: \begin{ayah}مِنْ قَبْلِ أَنْ يَأْتِيَ يَوْمٌ لَا بَيْعٌ فِيهِ وَلَا خُلَّةٌ وَلَا شَفَاعَةٌ\end{ayah}، فذلك يوم انقطع فيه زمن الاكتساب والعمل، فلا يشتري العبد فيه لنفسه منزلة بمال، ولا ينفعه خليل كما كان في الدنيا، ولا تقوم الشفاعة للكفار كما كانت تقوم لبعض الناس بينهم في الدنيا بأسباب القرابة والمودة.

ونبّه الشيخ إلى أن نفي الشفاعة هنا عام اللفظ، لكنه مخصوص المعنى، لأن الشفاعة الثابتة للمؤمنين جاءت بها نصوص أخرى. فليست هذه الآية نفياً مطلقاً لكل شفاعة، وإنما نفيٌ للشفاعة التي يتوهمها أهل الكفر أو من لا تنفعه شفاعة الشافعين.

\begin{fawaid}[أول النفقة أداء الحقوق قبل التوسع في التطوع]
من الخطأ أن يتوسع العبد في الصدقات المستحبة ويغفل عن الحقوق الواجبة في ماله؛ لأن أول ما يدخل في الأمر بالإنفاق هو ما أوجبه الله عليه لا ما زاد عليه فقط.
\end{fawaid}

\section{آية الكرسي: تنزيه وكمال وافتقار الخلق إليه}
وقف الشيخ مع آية الكرسي وقفات كثيرة، وبيّن أن \rawi{الطبري} فيها يقرر كمال حياة الله وقيوميته، ونفي السنة والنوم عنه، وسعة ملكه، وإحاطة علمه، وأنه لا يؤوده حفظ السماوات والأرض.

ومن ألطف ما نبه عليه الشيخ هنا أن قوله تعالى: \begin{ayah}وَلَا يَئُودُهُ حِفْظُهُمَا\end{ayah} فيه بيان كمال الحفظ مع كمال السهولة في حق الله؛ فالعبد قد يحفظ شيئاً مع تعب ومشقة، أما الله سبحانه فيحفظ كل شيء ولا يثقله ذلك ولا يشق عليه.

كما أشار إلى أن \rawi{الطبري} في بعض مواضع الصفات قد يحكي الأقوال ولا يقطع بواحد منها إذا رأى أن بعضها يتكامل ولا يتضاد، أو إذا لم تتبين له حجة فاصلة على تعيين بعض التفصيل.

\begin{fawaid}[كمال الفعل الإلهي يظهر في وقوعه بلا كلفة ولا عجز]
قد يقدر المخلوق على بعض الأفعال لكن مع تعب ومشقة ونقص، أما الله سبحانه فحفظه وتدبيره وعلمه وقدرته كلها كاملة لا يلحقها عجز ولا إعياء ولا ثقل.
\end{fawaid}

\section{لا إكراه في الدين: العموم والخصوص والنسخ}
عند قوله تعالى: \begin{ayah}لَا إِكْرَاهَ فِي الدِّينِ\end{ayah}، ذكر \rawi{الطبري} الأقوال الواردة في سبب النزول، ثم اختار أن الآية ليست منسوخة، ولكنها نزلت في خاص من الناس: ممن يقبل منهم بذل الجزية ويقرون على دينهم الباطل مع خضوعهم لحكم الإسلام.

وبيّن الشيخ أن هذا الموضع من أنفع المواضع في فهم الفرق بين أمرين:
\begin{itemize}
    \item أن يكون اللفظ عاماً ثم ينسخ بعض حكمه.
    \item وأن يكون ظاهر اللفظ عاماً لكن المراد به الخصوص من أول الأمر.
\end{itemize}

فالطبري هنا يرى أن الآية من الباب الثاني، لا من باب النسخ. ولذلك استثمر هذا الموضع في تقرير قاعدة أصولية جميلة: أن ما كان ظاهره العموم وباطنه الخصوص لا يدخل في باب الناسخ والمنسوخ.

\begin{fawaid}[ليس كل عموم ظاهري يقال فيه بالنسخ]
إذا دل الدليل على أن المراد بالآية من أول الأمر خاص، فليس هذا من باب النسخ، بل من باب العموم الذي أريد به الخصوص. وهذه قاعدة مهمة جداً في قراءة كتب التفسير والأصول.
\end{fawaid}

\section{قد تبين الرشد من الغي}
فسر \rawi{الطبري} الرشد بإصابة الحق، والغي بمجاوزته إلى الباطل، وبيّن أن المعنى: قد اتضح الحق من الباطل، واستبان طريق الهدى من طريق الضلال، فلا معنى لإكراه من ظهرت له الحجة وكان من أهل الإقرار على دينه إلى حين.

وأبرز الشيخ أن من أعظم معاني هذه الآية أن دين الله في نفسه مبني على الحجة والبيان، لا على الإلجاء القسري الذي يقوم مقام البرهان. فالحق ظاهر، والرشد بيّن، ومن أعرض بعد البيان فإلى الله أمره.

\section{الطاغوت والمعنى الجامع للاختلاف}
في قوله تعالى: \begin{ayah}فَمَنْ يَكْفُرْ بِالطَّاغُوتِ\end{ayah}، ذكر \rawi{الطبري} أقوالاً في معنى الطاغوت: فقيل هو الشيطان، وقيل الساحر، وقيل الكاهن، ثم رجح المعنى الجامع فقال: هو كل ذي طغيان عُبد من دون الله أو اتبع في معصية الله.

وهذا مثال بديع على ما يفعله \rawi{الطبري} كثيراً: يرد الأقوال المختلفة إلى أصل كلي يجمعها، فتظهر حينئذ أنها ليست أقوالاً متضادة، بل أمثلة وصور لمعنى واحد.

\begin{fawaid}[أحياناً يكون الخلاف تفسيراً بالمثال لا تضاداً في المعنى]
إذا اختلفت عبارات السلف في تفسير لفظ واحد، فقد يكون الصواب جمعها تحت أصل كلي، لأن كل قول منها يمثل صورة من صوره لا معنى مستقلاً يناقض غيره.
\end{fawaid}

\section{العروة الوثقى ومعنى الثبات}
فسر \rawi{الطبري} العروة الوثقى بأنها مثل للإيمان الذي يتعلق به العبد تعلق المستمسك بأوثق عروة لا يخشى انقطاعها. ومن هنا جاء قوله تعالى: \begin{ayah}لَا انْفِصَامَ لَهَا\end{ayah}، أي: لا انكسار لها.

ووقف الشيخ عند أثر جميل روي في هذا الباب، وفيه أن حسن الخاتمة يكون بالثبات على أصل الإيمان والكفر بالطاغوت. فكلما ثبت هذا الأصل في القلب زاد أمن العبد من الانفصام والانتكاس.

\section{الله ولي الذين آمنوا}
شرح \rawi{الطبري} أن ولاية الله للمؤمنين معناها توليه لهم بالنصرة والهداية والتوفيق، وأن إخراجهم من الظلمات إلى النور هو إخراج من ظلمات الكفر والشك والجهل إلى نور الإيمان واليقين.

ثم ذكر في مقابلة ذلك أن الطاغوت أولياء الكافرين يخرجونهم من النور إلى الظلمات، وذكر احتمالين في هذا الإخراج:
\begin{itemize}
    \item أن يكون إخراجاً ممن كان عنده شيء من الحق ثم انتكس عنه.
    \item أو أن يكون صدّاً ومنعاً من الدخول في النور أصلاً، فعبر عن الحرمان بالإخراج على جهة التوسع العربي في الاستعمال.
\end{itemize}

وفي هذا من دقة \rawi{الطبري} ما يبين كيف يوازن بين الحقيقة اللغوية والسياق، وكيف يفسر التعبير القرآني على أوسع ما يحتمله اللسان العربي الصحيح.

\begin{fawaid}[قد يعبر القرآن بالإخراج عن مجرد الصد والمنع]
من دقائق العربية في التفسير أن الإخراج قد يستعمل في منع الشيء وصرف العبد عنه، وإن لم يكن داخلاً فيه من قبل دخولاً مستقراً، وهذا يوسع فهم بعض التعابير القرآنية.
\end{fawaid}

\separator

\section{بداية محاجة إبراهيم عليه السلام}
ختم المجلس ببداية قوله تعالى: \begin{ayah}أَلَمْ تَرَ إِلَى الَّذِي حَاجَّ إِبْرَاهِيمَ فِي رَبِّهِ\end{ayah}، ثم دخل في قوله: \begin{ayah}وَإِذْ قَالَ إِبْرَاهِيمُ رَبِّ أَرِنِي كَيْفَ تُحْيِي الْمَوْتَى\end{ayah}.

وبيّن الشيخ أن \rawi{الطبري} ذكر عدة أقوال في سبب سؤال إبراهيم، ثم رجح أن ذلك كان لعارض ألقي في نفسه لا على وجه الشك المستقر، لكن الشيخ نبه إلى أن الحديث: \begin{hadith}نحن أحق بالشك من إبراهيم\end{hadith} ليس مراده إثبات الشك لإبراهيم، بل نفيه عنه من باب أولى؛ أي: إذا كنا نحن لا نشك فإبراهيم أولى بذلك.

ومن الفوائد النفيسة هنا أن الأنبياء إذا وقع منهم ما يقتضي استغفاراً أو توبة جاء ذلك صريحاً في القرآن، فلما لم يأت هنا ما يدل على توبة من شك أو استغفار منه، كان هذا من أقوى ما يدل على أن سؤاله لم يكن عن شك في قدرة الله، بل عن طلب مزيد طمأنينة وعيان.

\begin{fawaid}[طلب زيادة الطمأنينة ليس مناقضاً لأصل اليقين]
قد يكون العبد مؤمناً بالحق علماً وخبراً، ثم يطلب من الله مزيد كشف وطمأنينة وعيان، وهذا لا ينافي أصل الإيمان، بل هو من كمال التعلق بالله ومحبة ازدياد اليقين.
\end{fawaid}

\section{فائدة في قوله: (والله لا يهدي القوم الظالمين)}
نبّه الشيخ في أثناء هذا الموضع إلى معنى دقيق في قوله تعالى: \begin{ayah}وَاللَّهُ لَا يَهْدِي الْقَوْمَ الظَّالِمِينَ\end{ayah}، وهو أن نفي الهداية هنا لا يقتصر على أصل التوفيق العام، بل يدخل فيه أنه لا يهديهم إلى حجة صحيحة يدحضون بها أهل الحق عند الخصومة؛ لأن حجج الباطل داحضة من أصلها، وإنما قد يروج بعضها زمناً بما يقع من الشبهة والتمويه لا بقوة البرهان.

وهذا من المعاني التي تطمئن قلب المؤمن والداعية معاً: فصاحب الباطل قد يكثر الجدل، لكنه لا يرزق حجة مستقيمة يثبت بها على التحقيق، وأما أهل الحق فمعهم من الله من النور والبيان ما يثبتهم إذا أخلصوا وصدقوا.

\section{(أو كالذي مر على قرية): المقصود من القصة}
لما دخل \rawi{الطبري} إلى قوله تعالى: \begin{ayah}أَوْ كَالَّذِي مَرَّ عَلَى قَرْيَةٍ\end{ayah}، ذكر الخلاف في تعيين الرجل: أهو عزير أم إرميا أم غيرهما، ثم قرر أصله المنهجي المكرر: أن هذا التعيين لا تقوم عليه حجة يجب التسليم لها، ولا تتوقف عليه العبرة، فلا حاجة إلى التكلف في الجزم بما لم يأت به بيان صحيح.

ثم أوضح الشيخ أن \rawi{الطبري} لم يقف عند هذا الحد، بل صرح بمقصد سوق القصة في السورة: إقامة الحجة على منكري الإحياء بعد الموت من مشركي العرب، وتثبيت الحجة كذلك على يهود المدينة بإخبار النبي صلى الله عليه وسلم بما لا سبيل له إلى معرفته من قبل نفسه ولا من قبل قومه الأميين. فالقصة عنده ليست مجرد مشهد عجيب في إحياء الموتى، بل برهان مركب: برهان على البعث، وبرهان على الرسالة.

\begin{fawaid}[إذا تم المقصود من القصة لم يتوقف على تعيين أسمائها]
كثير من القصص القرآني يساق للعبرة وإقامة الحجة، لا للتوثيق التاريخي المجرد. فإذا ثبت هذا المقصود، لم يكن الاشتغال بتعيين الشخص أو الموضع أو الزمان أولى من تدبر الحجة التي سيقت القصة لأجلها.
\end{fawaid}

\separator

خلاصة هذا اللقاء أن \rawi{الطبري} جمع فيه بين تقرير أصول عقدية عظيمة في تفضيل الرسل، والإرادة، والشفاعة، والصفات، وبين قواعد أصولية دقيقة في باب العموم والخصوص والنسخ، مع أمثلة بديعة على رد الاختلاف إلى المعنى الجامع، وفتح باب لطيف في فهم سؤال إبراهيم عليه السلام على وجه يليق بمقام النبوة واليقين.
